% displayname : Applications des champs magnétostatiques
\documentclass[12pt,fancy]{cours}
\Chapit{5}{4}
\begin{document}
\maketitre{
Nous poursuivons l'étude de la magnétostatique avec l'entité magnétique la plus élémentaire: le dipôle. En effet l'équation de Maxwell-Thomson interdit l'existence de monopôle magnétique, le dipôle est donc l'entité la plus simple à étudier. Nous allons voir qu'on associe aux spires de courant un moment dipolaire magnétique. L'approximation dipolaire permet alors de déterminer le champ magnétique créé par ce dipôle à grande distance. Cette étude conduira à une description qualitative du fonctionnement d'un aimant et des processus d'aimantation de la matière.
}
{
\item Cours de Champ magnétique et induction  2 - 0 (induction)
\item Documents de cours Champ magnétique et induction 2 - 1 et 2
\item TD Champ magnétique et induction 2
}
%%%%%%%%%%%%%%%%%%%%%%%%%%%%%%%%%%%%%%


% \subsection{Exemple de la spire de courant}

% \begin{definition}[Spire de courant]\label{def:}
%   Une \imp{spire de courant} est un fil circulaire fermé parcouru par un courant électrique.
% \end{definition}
% On supposera la spire de forme circulaire. 
% La spire de courant est la structure électrique fondamentale en magnétostatique, de même que la charge ponctuelle pour l'électrostatique.
% L'allure du champ magnétostatique créé par une spire de courant (Fig.~\ref{fig:cartes_de_champ}) est à connaître. On vérifie aisément que les propriétés de symétrie sont satisfaites: d'une part le plan horizontal contenant la spire (perpendiculaire au plan de la feuille) est un plan de symétrie de la distribution de courants, donc $\B$ est orthogonal à ce plan; d'autre part le plan vertical perpendiculaire au plan de la feuille coupant la spire en deux est un plan d'antisymétrie de la distribution de courants, donc $\B$ est contenu dans ce plan. Ces deux propriétés sont bien vérifiées sur la Fig.~\ref{fig:cartes_de_champ}.
% \begin{figure}[H]
% \centering
% \includegraphics[width=0.6\textwidth]{spire_courant_B.pdf}
% \caption{Allure du champ magnétostatique créé par une spire de courant (en rouge) parcourue par un courant uniforme. La figure est à symétrie de révolution autour de l'axe vertical passant par le centre de la spire.}
% \label{fig:cartes_de_champ}
% \end{figure}



\section{Le solénoïde idéal}

%
%\paragraph{Et si le fil devient un cylindre infini de rayon $a$?}
%Les invariances et symétries sont identiques, et si la densité volumique de courant est uniforme alors $\vv{j}=\frac{I}{\pi a^2}\ez$. Le contour d'Ampère est le même, un cercle orienté dans le sens de $\vv{j}$ avec la règle de la main droite, de rayon $r$ centré sur $O$ et passant par $M$. Pour calculer le courant enlacé, il faut distinguer deux cas. Dans chaque cas, la surface $S(\Gamma)$ qui s'appuie sur le contour $\Gamma$ est un cylindre d'axe $(Oz)$ et de rayon $r$, et puisque $\vv{j}$ est selon $\ez$ dans le calcul du flux seule importe la surface $\d S=r\d r\d \theta$ de normale $\ez$.
%\begin{enumerate}
%\item Si $r<a$, alors $I_{\rm{enlacé}}=\iint_{S(\Gamma)}\vv{j}\cdot\vv{\d S}=j\pi r^2=\frac{I r^2}{a^2}$. Alors le théorème d'Ampère donne
%\begin{equation}
%\vv{B}(r<a)=\frac{\mu_0 I r}{2\pi a^2}\vv{e_\theta}.
%\end{equation}
%\item Si $r>a$, alors $I_{\rm{enlacé}}=\iint_{S(\Gamma)}\vv{j}\cdot\vv{\d S}=j\pi a^2=I$. Alors le théorème d'Ampère donne
%\begin{equation}
%\vv{B}(r>a)=\frac{\mu_0 I }{2\pi r}\vv{e_\theta}.
%\end{equation}
%\end{enumerate}
%
%\begin{remarque}
%Si on ne tient pas compte du blindage du câble électrique, le champ créé à une distance d'environ \SI{1}{cm} pour un courant de \SI{1}{A} est de \SI{2e-5}{T}. C'est une valeur très faible, de l'ordre de grandeur du champ magnétique terrestre en France (ce champ est créé par les courants électriques de convection à l'intérieur du noyau externe terrestre qui contient de grandes quantités de fer liquide).
%\end{remarque}

%\subsection{Solénoïde infiniment long}
\subsection{Modèle et champ magnétostatique}

\parag{Le solénoïde}
Nous allons démontrer à partir des équations de Maxwell l'expression admise en première année du champ créé par un solénoïde. Rappelons qu’un \imp{solénoïde} est une bobine composé de $N$ spires jointives circulaires enroulées sur un cylindre de longueur $\ell$ et de rayon $a$, voir \Fig{sole}. Les solénoïdes permettent de réaliser des bobines à induction et des électroaimants. On caractérise un solénoïde par le nombre de spires par unité de longueur $n$ tel que $N=n\ell$. . On se place dans la limite du solénoïde infiniment long selon l'axe $(Oz)$, soit $L\gg a$, de sorte qu'à l'intérieur du solénoïde on puisse négliger les effets de bords.

\begin{definition}[Solénoïde idéal]\label{def:}
  Um \imp{solénoïde idéal} est un enroulement infiniment long de spires parcourues par un courant $I$.
\end{definition}

\begin{figure}[h!]
\centering
\subfigure[]{\label{fig:}
\includegraphics[width=0.45\textwidth]{solenoide}
}\hfill
\subfigure[]{\label{fig:}
\includegraphics[width=0.45\textwidth]{sole_field}
}
\caption{(a) Un solénoïde est un enroulement de longueur $\ell$ de $N=n\ell$ spires circulaires parcourues par un courant d'intensité $I$. (b) Carte de champ magnétostatique d’un solénoïde réel.} 
\label{fig:sole}
\end{figure}

\parag{Invariances et symétries} On commence toujours par cette étape.
\begin{liste}
  \item \udl{Symétrie} : le plan $\Pi_{\rm S}=(M,\er,\et)$ normal à $\ez$ est un plan de symétrie de la distribution de courants, et donc le champ est orthogonal à ce plan, soit $\vv{B}=B_z(M)\ez$.
  \item \udl{Invariance} : Le solénoïde  possède une symétrie cylindrique. La distribution de courants est donc invariante par rotation et translation selon l'axe $(Oz)$. Donc le champ ne dépend pas de $\theta$ ni de~$z$, soit 
  \begin{equation}
    \vv{B}(M) = B_z(r) \ez.
  \end{equation}
  
\end{liste}
% Finalement on en conclut qu'en tout point $M(r,\theta,z)$ on a $\vv{B}(M)=B_z(r)\ez$.

\parag{Théorème d’Ampère}

 Utilisons alors le théorème d'Ampère pour calculer le champ. Puisque le champ est selon $\ez$, il faut un contour qui passe par $M$ et dont une portion est parallèle à $(Oz)$. Le plus simple est ainsi de choisir un rectangle (Fig.~\ref{fig:ampère_solénoide}). Les portions verticales du contour sont orthogonales au champ ($\vv{B}\cdot\vv{\d l}_{\rm{vert}}=0$) et ne contribuent donc pas à la circulation. Nous distinguons trois situations suivant que la position du contour d'Ampère. Le calcul de la circulation est commun aux trois cas.

\begin{figure}[h!]
\centering
\subfigure[]{\label{fig:}
\includegraphics[width=0.3\textwidth]{solenoide_ampère_a.pdf}
}\hfill
\subfigure[]{\label{fig:}
\includegraphics[width=0.3\textwidth]{solenoide_ampère_b.pdf}
}\hfill
\subfigure[]{\label{fig:}
\includegraphics[width=0.3\textwidth]{solenoide_ampère_c.pdf}
}
\caption{Choix du contour d'Ampère pour le solénoïde supposé infini, dans le cas où le point $M$ est à l'intérieur du solénoïde $(r<a$) et dans le cas où il est à l'extérieur ($r>a$).}
\label{fig:ampère_solénoide}
\end{figure}

\begin{liste}
\item Pour $r<a$, le théorème d'Ampère donne
\begin{equation}
\int_{\Gamma}\vv{B}\cdot\vv{\d l}=B(r)h-B_{\rm{axe}}h=\mu_0 I_{\rm{enla}}=0
\end{equation}
car le contour est situé à l'intérieur de la bobine et il n'enlace donc aucun courant. Le champ dans le solénoïde est donc uniforme (indépendant de $r$) !
%et il vaut la valeur du champ sur l'axe (en $r=0$):
\begin{equation}
\vv{B}(r<a)=B_{\rm{axe}}\,\ez.
\end{equation}

\item
Pour $r>a$, le théorème d'Ampère donne
\begin{equation}
\int_{\Gamma}\vv{B}\cdot\vv{\d l}=B(r)h-B_{\rm{axe}}h=\mu_0 I_{\rm{enla}}=-\mu_0 N\frac{h}{L}I = -\mu_0 nhI,
\end{equation}
où $N\frac{h}{L}$ est le nombre de spires enlacées par le contour $\Gamma$.

\begin{remarque}
Attention au signe, ici $I_{\rm{enla}}<0$ car le contour est orienté dans le sens direct (inverse des aiguilles d’une montre) mais $I$ sort vers l’arrière du plan.
\end{remarque}
\end{liste}
\begin{equation}
\vv{B}(r>a)=(B_{\rm{axe}}-\mu_0 n I)\,\ez.
\end{equation}
Le champ est aussi uniforme en dehors du solénoïde. En particulier il est égal à sa valeur en $r=\infty$, qui est nulle. On peut donc remplacer $\Vec{B}(r>a)$ par $\Vec{0}$, ce qui donne $B_{\rm axe} = \mu_0 n I$. En conclusion,
\begin{theorem}[Champ magnétostatique d’un solénoïde infini]\label{form:champ_sol}
\begin{equation*}
\vv{B}(r<a)=\mu_0 n I\,\ez \text{ (à l’intérieur)}, \qquad
\vv{B}(r>a)=\Vec{0} \text{ (à l’extérieur)}.
\end{equation*}
\end{theorem}

\begin{figure}[h!]
\centering

\label{fig:}
\end{figure}
%%%%%%%%%%%%%


\subsection{Inductance et énergie stockée}
\parag{Inductance} Soit un circuit parcouru par un courant d'intensité $I$ uniforme. Notons $\Phi$ le flux magnétique propre, c'est-à-dire le flux du champ magnétique créé par le courant d'intensité $I$ à travers le circuit-même. Le flux magnétique à travers une spire circulaire $\Surf$ de rayon $a$ vaut 
\begin{equation}
  \Phi_s = \iint_{\Surf} \vv{B}\cdot\vv{\d S} = B\times \pi a^2 = \mu_0 n I \pi a^2.
\end{equation}
Le flux à travers l'ensemble du solénoïde, qui comprend $N=n\ell$ spires, vaut donc
\begin{equation}
  \Phi = N \Phi_s = \mu_0 n^2 \pi a^2 \ell I.
\end{equation} 
On remarque que le flux est proportionnelle à l'intensité du courant $I$.
\begin{definition}[Inductance propre]
L'\imp{inductance propre} est la grandeur $\mc{L}>0$ (en henry H) reliée au flux magnétique propre $\Phi$ par
% \footnote{D'habitude notée $L$ mais ici $L$ désigne déjà la longueur du solénoïde!}
\begin{equation*}
\Phi=\mc{L} I.
\end{equation*}
\end{definition}
\noindent
% Pour un solénoïde de rayon $a$ comportant $N$ spires de densité linéique $n=N/L$, le flux vaut
% \begin{equation}
% \begin{split}
% \Phi=\iint_{S}\vv{B}\cdot\vv{\d S}= B\times N\times\pi a^2=\left(\mu_0 N n\pi a^2 \right)\times I,
% \end{split}
% \end{equation}
% d’après la \Formu{champ_sol}.
 L’inductance d’un solénoïde idéal vaut donc
% \begin{equation}
\eqbox[inductance]{
  \mc{L} = \mu_0 N n\pi a^2.
  }
% \end{equation}

\parag{Énergie stockée}
%Pour terminer, exprimons l'énergie stockée dans la bobine.
La puissance électrique consommé par un dipôle s'écrit $P=UI$. Vous avez vu en première année que la caractéristique de la bobine est\footnote{C'est une conséquence de la loi de Faraday (obtenue à partir de l'équation de Maxwell-Faraday en régime variable), que nous verrons bientôt.}
\begin{equation}
U=\mc{L}\frac{\d I}{\d t}.
\end{equation}
Comme $\mc{L}=\cte$, on peut écrire $P=\frac{\d}{\d t}\left(\frac{1}{2}\mc{L}I^2\right) = \dd{\mc{E}_{\mc{L}}}{t}$. On peut alors identifier l'énergie stockée dans la bobine à
\begin{equation}
\mc{E}_{\mc{L}}=\frac{1}{2}\mc{L}I^2.
\end{equation}
Or $\mc{L}=\mu_0 Nn\pi a^2$ et $I=\frac{B}{\mu_0 n}$, donc l'énergie stockée dans un solénoïde se met sous la forme
\begin{equation}
\mc{E}_{\mc{L}}=\frac{B^2}{2\mu_0}\times (\pi a^2 L).
\end{equation}
Le terme $V = \pi a^2 L$ désigne le volume intérieur de la bobine. Ainsi, le facteur $\frac{B^2}{2\mu_0}$ représente la densité volumique d'énergie magnétique stockée dans la bobine. Cette expression est en fait très générale.
\begin{theorem}[Densité volumique d'énergie magnétique]
\begin{equation*}
e_{\rm m} = \frac{B^2}{2\mu_0}
\end{equation*}
\end{theorem}

\section{Le dipôle magnétostatique}

Le dipôle magnétique est l’équivalent du dipôle électrostatique. Dans cette section, nous modélisons un dipôle magnétique par une boucle de courant, puis déterminons l’expression du champ magnétostatique dipolaire.
% Enfin, nous étudions l’action d’un champ magnétique extérieur sur un dipôle magnétique.


\subsection{Modélisation par une spire de courant}
\label{sec:dipole}


\begin{definition}[Spire de courant]\label{def:}
  Une \imp{spire de courant} est un fil circulaire fermé parcouru par un courant électrique.
\end{definition}
On supposera la spire de forme circulaire. 
La spire de courant est la structure électrique fondamentale en magnétostatique, de même que la charge ponctuelle pour l'électrostatique.
Soit une spire de courant orientée $\Gamma$ parcourue par un courant uniforme d'intensité $I$. Soit $\vv{n}$ la normale à la section $\Surf$ de la spire qui s'appuie sur le contour orienté $\Gamma$ obtenu par la règle de la main droite. On note $d$ le diamètre de la spire et $S$ sa surface. On associe à cette spire un \imp{moment (dipolaire) magnétique} par la relation suivante\footnote{Le plus souvent, le moment dipolaire électrique est simplement appelé moment dipolaire, et le moment dipolaire magnétique simplement moment magnétique.}.
\begin{definition}[Moment magnétique]
Le \imp{moment magnétique} d’une spire de courant (en $\SI{}{A.m^2}$) est défini par
\begin{equation*}
 \M=IS\,\vv{n}.
 \end{equation*}
\end{definition}
% \noindent Pour une spire circulaire, comme sur la \Fig{spire}, $S=\pi R^2$ où $R$ est le rayon de la spire.
%%%%%%%%%%%%%%%%%
\begin{remarque}
  On peut généraliser à un contour $\Gamma$ de forme quelconque parcouru par un courant, par la relation 
  \begin{equation}
  \M=I\oint_{\Gamma}\vv{r}\pv\vv{\d l}.
  \end{equation}
\end{remarque}

 \begin{figure}[H]
 \centering
 \includegraphics[width=0.26\textwidth]{spire_surface.pdf}
 \caption{Une spire de courant génère un moment magnétique $\M$ normal à la spire.}
 \label{fig:spire}
 \end{figure}
%%%%%%%%%%%%%%%

On peut appliquer le moment magnétique à d'autres distributions similaires à la boucle de courant. La spire de courant n'est qu'un modèle particulier d'une distribution magnétique plus générale appelée dipôle magnétostatique.

\begin{definition}[Dipôle magnétostatique]\label{def:}
  On appelle \imp{dipôle magnétostatique} toute distribution de courant stationnaire de taille  $d$ négligeable, i.e. étudié à grande distance $r \gg d$ (approximation dipolaire).
\end{definition}

On peut associer à tout dipôle magnétostatique un moment magnétique $\M$.

\subsection{Analogie avec le dipôle électrostatique}

%%%%%%%%%%%%%%%%%%%%%
\begin{figure}[t!]
\centering
\includegraphics[width=0.98\textwidth]{dipoles_champ_comparaison.pdf}
\caption{À gauche, le champ magnétostatique créé par une spire de courant, à droite le champ électrostatique créé par un doublet de charges (positive en rouge, négative en bleu). Hormis dans la région proche des sources représentée en orange, les deux cartes de champ sont très proches.
% Ceci justifie d'une part que loin de la spire de courant, on peut assimiler la spire à un dipôle magnétique et d'autre part que l'on peut prendre une expression analogue pour le dipôle magnétostatique à celle obtenue pour le champ du dipôle électrostatique.
}
\label{fig:spire_et_doublet_de_charges}
\end{figure}
%%%%%%%%%%%%%%%%

\begin{remarque}
  Contrairement à l'électorastique, il n'existe pas de charge magnétique. On ne peut donc pas modéliser un dipôle magnétostatique par un doublet de charges, comme pour le dipôle életcrostatique. Ces deux types de dipôles présentent malgré tout beaucoup de similarités.
\end{remarque}

Le champ magnétostatique créé par une boucle de courant est analogue au champ électorastique créé par un doublet de charges électriques. Sa carte de champ magnétique ressemble fortement à celle du champ électrostatique créé par un doublet, en particulier si on se place loin des sources (Fig.~\ref{fig:spire_et_doublet_de_charges}).
 Dans l'\imp{approximation dipolaire} $r \gg d$, la structure interne des sources n’a plus d’importance, et les cartes de champ sont identiques. On peut s’en convaincre en comparant les invariances et symétries d’une boucle de courant et d’un doublet de charge.

 \begin{liste}
 \item Le plan $\Pi_{x}$ de normale $\ex$ est un plan d’antisymétrie de la boucle de courant et de symétrie du doublet de charge : donc les cartes de $\Vec{B}$ et $\Vec{E}$ sont symétriques par rapport à $\Pi_x$.
  \item Le plan $\Pi_{y}$ de normale $\ey$ est un plan de symétrie de la boucle de courant et d’antisymétrie du doublet de charge : donc les cartes de $\Vec{B}$ et $\Vec{E}$ sont antisymétriques par rapport à $\Pi_y$.
 \end{liste}


%\begin{remarque}
%On peut se convaincre que loin des sources il est logique de retrouver la même allure du champ, car dans l'espace vide, les champs $\vv{E}$ et $\vv{B}$ obéissent aux mêmes équations: $\dive{\vv{E}}=0$ (en l'absence de charge), $\rot{\vv{E}}=0$ et $\dive{\vv{B}}=0$, $\rot{\vv{B}}=0$ (en l'absence de courant).
%\end{remarque}

%Autrement dit en des points suffisamment éloignés d'une boucle de courant, le champ magnétique créé tend vers celui d'un dipôle magnétique de moment $\M$, de sorte que
On peut retenir de façon schématique que
\begin{prop}\label{prop:}
  Dans l'approximation dipolaire $r \gg d$, le champ magnétique créé par une boucle de courant est analogue au champ électrostatique créé par un doublet de charges, soit
\begin{equation*}
\text{dipôle magnétostatique } (spire~de~courant)\quad\Leftrightarrow\quad \text{dipôle électrostatique } (doublet~de~charges)
\end{equation*}
\end{prop}


%Les sources du magnétisme dans la matière sont des dipôles magnétiques, qui seront toujours assimilables à de petites boucles de courant.

\subsection{Expression du champ dipolaire}

En raisonnant par analogie avec le dipôle électrostatique, on remplace
\begin{equation}
\begin{cases}
\vv{E}\text{ par }\vv{B},\\
\vv{p}\text{ par }\M,\\
\frac{1}{\varepsilon_0}\text{ par }\mu_0.
\end{cases}
\end{equation}
En posant $\M=\mu\,\ez=\M\,(\cos\theta\,\vv{e_r}-\sin\theta\,\vv{e_\theta})$, le champ dipolaire magnétique s'écrit donc
\begin{equation}
\label{eq:champ_E_dipolaire}
\vv{B}(M)=\frac{\mu_0 \M}{4\pi r^3} \left(2\cos\theta\,\vv{e_r}+\sin\theta\,\vv{e_\theta}\right).
\end{equation}
Le champ magnétostatique se met aussi sous une forme indépendante du système de coordonnées.
\begin{theorem}[Champ magnétostatique dipolaire]
\begin{equation}
\B_{\rm dip}(M)=\frac{\mu_0}{4\pi}\frac{3(\M\cdot\vv{r})\,\vv{r}-r^2\M}{r^5}.
\end{equation}
\end{theorem}
On rappelle que $\frac{\mu_0}{4\pi} = \SI{e-7}{H.m^{-1}}$.




\subsection{Actions d'un champ extérieur sur un dipôle}
On s’intéresse à l’action d’un champ magnétique extérieur $\B_{\rm ext}$ sur un dipôle magnétique $\M$.

\parag{Résultante des forces}
\label{subsec:résultante}
La spire de courant constitue un circuit longiligne fermé. La force qui s'exerce sur une portion de fil est la force de Laplace.

\begin{theorem}[Force de Laplace]
  Une portion $\d l$  de fil parcourue par une intensité $I$ dans un champ magnétique $\vv{B}$ subit la force
  \begin{equation*}
\d \vv{F}_{\rm{L}}= I\vv{\d l}\pv \vv{B}.
\end{equation*}
\end{theorem}


\udl{Champ uniforme}.
Supposons pour commencer que le champ extérieur $\B_{\rm ext}=\Vec{B}_0$ est uniforme. 
Alors la résultante des forces de Laplace qui s'exercent sur la spire s'écrit
\begin{equation}
\vv{F}_{\rm{L}}=\oint_{\Gamma}I\vv{\d l}\pv \vv{B_0}=I\underbrace{\left(\oint_{\Gamma}\vv{\d l}\right)}_{=\Vec{0}}\pv\,\vv{B_0}=\Vec{0}.
\end{equation}
\begin{remarque}
  On peut se convaincre de la nullité de l'intégrale dans le cas d'une boucle circulaire de rayon $R$. On a alors 
  $$
\oint_{\Gamma}\vv{\d l}=\int_0^{2\pi}R\left(-\sin\theta\,\vv{e_x}+\cos\theta\,\vv{e_y}\right)\d \theta= R \left[-\cos\theta\right]_0^{2\pi} \ex + R \left[\sin\theta\right]_0^{2\pi} \ey = \Vec{0}.
  $$
  Ce résultat est valable sur tout contour fermé.

\end{remarque}
\begin{prop}\label{prop:}
Un champ magnétostatique \imp{uniforme} exerce une résultante des forces \imp{nulle} sur un dipôle magnétique.
\end{prop}
\noindent
Nous avions établi un résultat analogue pour le dipôle électrostatique.

\udl{Champ non uniforme}. De la même manière, si le champ extérieur $\B_{\rm ext}$ varie sur une distance typique petite ou similaire à la taille~$d$ du dipôle magnétique, la résultante $\vv{F}$ n’est pas nulle. Par analogie avec le dipôle électrostatique, on remplace $\Vec{E}$ par $\Vec{B}$ et $\Vec{p}$ par $\M$.
%, et donc la force  magnétique s’exerçant sur un dipôle magnétique s’écrit
\begin{theorem}[Force subie par un moment magnétique]
\begin{equation*}
\vv{F}=\left(\M\cdot\grad{}\right)\B_{\rm ext}
\end{equation*}
\end{theorem}
\noindent

% \begin{remarque}
% On rappelle que le moment magnétique $\Vec{M}$ doit être pris constant sous l'opérateur gradient (autrement dit $\M$ est fixé lors de la dérivation). Cette précision a son importance lorsque $\Vec{M}$ dépend du champ magnétique $\Vec{B}(O)$ au point $O$ où se situe le moment magnétique. Dans ce cas, $\Vec{B}(O)$ est bien un vecteur constant, et ne doit pas être assimilé à $\Vec{B}(M)$ dans la dérivation !
% \end{remarque}

\parag{Résultante des moments}
On procède par analogie avec le dipôle électrostatique.
\begin{theorem}[Moment subi par un dipôle magnétique]\label{form:Gamma_dip}
\begin{equation*}
\vv{\Gamma}=\M\pv\B_{\rm ext}
\end{equation*}
\end{theorem}
\noindent Soit $\theta=(\M,\B_{\rm ext})$ l'angle entre le moment magnétique et le champ magnétique. Le moment de norme $\Gamma=\mu B\sin\theta$ fait tourner le dipôle jusqu'à ce qu'il soit aligné avec le champ ($\theta=0$).


\begin{remarque}
Pour obtenir la \Formu{Gamma_dip}, on rappelle qu’au premier ordre le champ $\B_{\rm ext}$ est supposé uniforme sur la taille de la boucle de courant, de sorte que le moment $\vv{\Gamma}$ ne dépend pas du point par rapport auquel on le calcule. Ce moment tend à aligner le dipôle magnétique avec le champ.
\end{remarque}

\parag{Énergie potentielle d'interaction}
La force magnétique exercée par le champ extérieur sur une boucle de courant \imp{rigide} est conservative. La démonstration est similaire à celle du dipôle électrostatique.
\begin{theorem}[Énergie potentielle d'un moment magnétique]
\begin{equation*}
  \Epot = -\M\cdot\vv{B}_{\rm ext}.
  \end{equation*}
\end{theorem}
% \begin{demo}
%   Pour une spire de courant \imp{rigide}, on a 
%   \eqbox[]{
%   \M = \cte.
%   }
%  En utilisant $\rot{\vv{B}}=0$, on peut aussi écrire la force pour une spire indéformable sous la forme
% \begin{equation}\label{eq:Fmag}
% \vv{F}=\grad{(\M\cdot\vv{B})}.
% \end{equation}
% Cette force dérive d'une énergie potentielle telle que $\vv{F}=-\grad{\Epot}$ avec à une constante près
% \end{demo}
\noindent
Cette énergie est minimale lorsque $\M$ et $\vv{B}$ sont colinéaires et de même sens (équilibre stable). Elle est maximale lorsque $\M$ et $\vv{B}$ sont colinéaires et de sens opposé (équilibre instable). En effet, l'énergie potentielle s'écrit aussi $-\mu B\cos\theta$, et est minimale pour $\theta=0$ et maximale pour $\theta=\pi$.
% S'ils sont dans des sens opposés, l'équilibre du dipôle est instable (énergie potentielle maximale).

%\subsubsection{Bilan: mouvement d'un dipôle dans un champ extérieur}
%Le moment, de norme $\Gamma=\MB\sin\theta$ fait tourner le dipôle jusqu'à ce qu'il soit aligné avec le champ ($\theta=0$). Dans ce cas, la force exercée par le champ extérieur sur le dipôle s'écrit $\M\,\grad{(B\cos\theta)}=\M\,\grad{B}$, et elle est effectivement orientée vers les régions où le champ est le plus intense (régions où les lignes de champ sont resserrées sur une carte de champ).

\section{Magnétisme dans la matière}
Nous discutons de l'observation du magnétisme à l'échelle macroscopique en \Sec{aimants_permanents}, puis nous discutons de l'origine microscopique du magnétisme dans la matière en \Sec{orbital} et \Sec{spin}.
% À l'échelle microscopique, les effets quantiques sont indispensables pour expliquer le magnétisme. En particulier, comme nous le verrons dans les chapitres de mécanique quantique, parler de la trajectoire d'un électron n'a pas de sens.   Néanmoins, certains modèles classiques (c'est-à-dire non quantiques) permettent de rendre compte des résultats observés.
% %Entendons-nous bien: les modèles proposés ne rendent pas compte de la réalité des phénomènes à l'échelle microscopique, néanmoins le résultat final obtenu est en accord avec le résultat rigoureux que l'on peut établir en mécanique quantique.
% Dans la suite nous allons donc prétendre que \guill{tout se passe comme si}, puisque les résultats que nous allons établir sont cohérents avec les résultats rigoureux de la théorie quantique.


\subsection{Aimants permanents}
\label{sec:aimants_permanents}

\parag{Définitions}
Du point de vue de leur réponse magnétique, les matériaux se scindent en trois groupes.
\begin{enumerate}
  \item Les matériaux \imp{diamagnétiques} sont repoussés par un champ magnétique extérieur. Ils ne possèdent pas de moment magnétique permanent, mais développent un moment magnétique induit dans le sens opposé au champ extérieur, soit $\M \propto -\B_{\rm ext}$. 
  
  \hspace{1cm}\udl{Exemples} : le cuivre Cu, l'argent Ag, le bismuth Bi, l'eau H$_2$O, les matériaux organiques.
  \item Les matériaux \imp{paramagnétiques} sont attirés par un champ magnétique extérieur. Ils ne possèdent pas de moment magnétique permanent, mais développent un moment magnétique induit dans le même sens que le champ extérieur, soit $\M \propto \B_{\rm ext}$.
  
  \hspace{1cm}\udl{Exemples} : l'aluminium Al, le tungstène W, le platine Pt, l'oxygène O$_2$.
  \item Les matériaux \imp{ferromagnétiques} sont attirés fortement par un champ magnétique extérieur, et peuvent posséder un moment magnétique permanent, en l'absence de champ magnétique extérieur.
  
  \hspace{1cm}\udl{Exemples} : le fer Fe, le cobalt Co, le nickel Ni, et leurs alliages : magnétite Fe$_3$O$_4$, ferrite, etc.
\end{enumerate} 

% , comme la magnétite, constituée essentiellement d'oyxde de fer Fe$_3$O$_4$, peuvent avoir un moment magnétique non nul en l'absence de champ magnétique extérieur. Ces matériaux sont qualifiés de ferromagnétiques. Un matériau magnétique qui possède un moment magnétique non nul est appelé \imp{aimant permanent}. Pour devenir un aimant, un matériau ferromagnétique doit avoir ces moments magnétiques atomiques alignés en moyenne dans le même sens.
\begin{definition}[Aimant permanent]\label{def:}
Un \imp{aimant permanent} est de moment magnétique $\M$ non nul en l'\imp{absence} de champ extérieur. 

Le moment d'un aimant est orienté conventionnellement du pôle \imp{sud} vers le pôle \imp{nord} de l'aimant.
\end{definition}


\begin{figure}[H]
  \centering
  % \includegraphics[height=0.2\textwidth]{neodym.jpeg}\hspace{0.5cm}
  % \subfigure{\includegraphics[height=0.2\textwidth]{weiss.png}}
  \subfigure{\includegraphics[height=0.4\textwidth]{magnet_field.pdf}}
  \subfigure{\includegraphics[height=0.4\textwidth]{champ_B_terre.png}}
  \caption{
    % (a) Un matériau ferromagnétique contient un grand nombre de domaines mésoscopiques de moments magnétiques bien définis (domaines de Weiss). Lorsque les moments magnétiques de ces domaines sont tous alignés, le matériau possède un moment magnétique total non nul : c'est un aimant permanent. 
    (a) Carte de champ magnétique créé par un aimant permanent. (b) Le champ magnétique terrestre est assimilable en bonne approximation à celui d'un dipôle magnétique.
  }
  \label{fig:weiss}
\end{figure}

Un milieu constituant un aimant permanent est caractérisé par son vecteur aimantation, qui correspond à la densité volumique de moment magnétique.
\begin{definition}[Aimantation]\label{def:}
L'\imp{aimantation} correspond au moment magnétique par unité de volume :
\begin{equation*}
\vv{M}=\frac{\d\M}{\d V} \quad \text{(en \SI{}{A.m^{-1}})}.
\end{equation*}
% où $\d \M$ désigne le moment magnétique d'un volume élémentaire $\d V$.
\end{definition}
\noindent
% L'aimantation $\vv{M}$ est donc une grandeur macroscopique intensive, exprimée en \SI{}{A.m^{-1}}.

\begin{table}[h!]
  \centering
  \begin{tabular}{|c|c|c|}
    \hline
    \textbf{Aimant} & \textbf{Moment magnétique $\mu$} & \textbf{Aimantation $M$} \\
    \hline
    Aimant au néodyme & $\sim \SI{1}{A.m^2}$ & $\sim \SI{e3}{A.m^{-1}}$ \\  
    Aimant de levage & $\sim \SI{e2}{A.m^2}$ & $\sim \SI{e5}{A.m^{-1}}$ \\
    Terre & $\sim \SI{e22}{A.m^2}$ & $\sim \SI{e5}{A.m^{-1}}$ \\ 
    Étoile à neutrons & $\sim \SI{e38}{A.m^2}$ & $\sim \SI{e11}{A.m^{-1}}$ \\
    \hline
  \end{tabular}
  \caption{Ordres de grandeur du moment magnétique $\mu$ et de l'aimantation $M$ pour différents aimants.}
  \label{tab:ordres_grandeur_aimants}
\end{table}

% \begin{exemple}{Aimantation maximale d'un aimant}
%   L'aimantation maximale $M_{\rm max}$ d'un milieu magnétique s'obtient quand tous les atomes sont parfaitement alignés. Chaque atome a un moment magnétique de l'ordre de $\mu_B$, de sorte que
%   \begin{equation}
%     M_{\rm max} \sim n \mu_B,
%   \end{equation}
%   où $n$ est le nombre d'atome par unité de volume.
%   Un milieu magnétique possède environ $n=\SI{e29}{}$ atomes par $m^3$. L'aimantation maximale vaut donc $M_{\rm max} \sim \SI{e6}{A.m^{-1}}$. Un aimant de volume $V=\SI{1}{cm^3}$ possède donc un moment magnétique de l'ordre de $\M = M_{\rm max} V \sim \SI{1}{A.m^2}$.

%   \end{exemple}

\parag{Force d'adhérence d'un aimant}
Les systèmes de verrous magétiques utilisent la force exercée entre deux aimants en contact pour assurer la fermeture d'un mécanisme. Comment relier cette force aux propriétés du matériau constituant l'aimant ?
\begin{theorem}[Force surfacique d'adhérence]\label{form:}
  La \imp{force surfacique d'adhérence} entre deux aimants en contact vaut
  \begin{equation*}
\frac{\d F}{\d S} \sim \frac{B^2}{\mu_0} \sim \mu_0 M^2.
\end{equation*}
\end{theorem}

\begin{demo}
Raisonnons par \udl{analyse dimensionnelle}. Soient deux aimants élémentaires de volume $\d V$, de moment magnétique $\d \M$, en contact sur une interface de surface $\d S$. Chaque aimant produit un champ magnétique dipolaire 
\begin{equation}
  B \sim \frac{\mu_0 \d \M}{4\pi r^3} \sim \frac{\mu_0 \d \M}{\d V} \sim \mu_0 M.
\end{equation}
Nous avons vu dans le chapitre précédent que la densité volumique d'énergie magnétique vaut
\begin{equation}
e_{\rm B} \sim  \frac{B^2}{\mu_0}.
\end{equation}
 Or, une énergie volumique est homogène à une force surfacique. On peut ainsi conclure qu'à un préfacteur numérique près, la force surface d'adhérence vaut
% \begin{equation}
\eqbox[]{
  \frac{\d F}{\d S} \sim \frac{B^2}{\mu_0} \sim \mu_0 M^2.
  }
% \end{equation}
\end{demo}

% \noindent En ODG, on trouve $\d F/\d S \sim \SI{e6}{Pa} = \SI{10}{bar}$.


\begin{exemple}{Aimant de levage}
On souhaite soulever une voiture de masse $m=\SI{1}{tonne}$ avec un gros aimant plat de section \SI{1}{m^2}. Déterminer le champ magnétique $B$ et l'aimantation $M$ de l'aimant pour réaliser cette tâche ?

\end{exemple}
\examplesol{
  \begin{quest}
    \item	
    D'après la relation précédente, la pression à exercée vaut 
    $$
    \frac{F}{S} = \frac{mg}{S} \sim \SI{e4}{Pa}.
    $$
    On en déduit le champ magnétique nécessaire pour soulever la voiture :
    % $$B\sim=\sqrt{\frac{\mu_0 mg}{S}},$$
    % soit 
    \eqbox{
    B \sim  \sqrt{\frac{\mu_0 F}{S}} \sim \SI{0,1}{T}
    }
    \item En en déduit l'aimantation de l'aimant :
    \eqbox{
    M \sim \frac{B}{\mu_0} \sim \SI{e5}{A.m^{-1}}.
    }
  \end{quest}
}

\subsection{Magnétisme orbital}
Nous discutons maintenant de l’origine \imp{microscopique} du magnétisme. Nous verrons que les moments magnétiques des atoms proviennent de deux contributions :
\begin{liste}
  \item le moment cinétique orbital des électrons autour du noyau atomique ;
  \item le spin, moment cinétique intrinsèque des électrons, que nous étudierons dans la section suivante.
\end{liste}


\parag{Modèle classique de l'atome d'hydrogène}
Les moments magnétiques existent à l'échelle atomique. En effet, les électrons sont des particules chargées en mouvement autour des noyaux atomiques. Ils créent donc des boucles de courant microscopiques, qui engendrent un moment magnétique, voir \Fig{modele_planetaire}.

Nous traitons l'atome d'hydrogène dans le \imp{modèle planétaire}, dont les hypothèses sont les suivantes. 
\begin{enumerate}[label=\fbox{H\arabic*}, leftmargin=2cm]
  \item Le proton est \imp{fixe} au centre de l'atome dans le référentiel du laboratoire ($\Rcal$). Cela revient à supposer que le proton est de masse $m_p \gg m_e$ très grande devant la masse de l'électron.
  \item L'électron décrit une orbite \imp{circulaire} de rayon $r$ autour du proton. Nous supposons que le mouvement de l'électron est donc bien décrit par la physique classique.
  \item La force exercée par le proton sur l'électron est de nature \imp{électrostatique}. Cela revient à supposer que la vitesse de l'électron  $v\ll c$ est petite devant la célérité de la lumière.
\end{enumerate}
\begin{figure}[h!]
\centering
\includegraphics[width=0.35\textwidth]{modele_planetaire.pdf}
\caption{Modèle planétaire de l'atome. L'électron en $\rm M$ décrit une orbite circulaire autour du noyau en $\rm O$.}
\label{fig:modele_planetaire}
\end{figure}
Appliquons la méthode du TMC pour analyser ce modèle.
\begin{enumerate}
  \item \udl{Référentiel} : référentiel du laboratoire $(\Rcal)$ dans lequel le proton est immobile.
  \item \udl{Système} : l'électron de masse $m_e$ et de charge $-e$.
  \item \udl{Moments} : la force électrostatique de Coulomb  $ \F_{\rm C}$
  % \begin{equation}
    % = -\frac{e^2}{4\pi \epz r^2}\vv{e_r},
  % \end{equation}
  est centrale, et donc de moment nul : 
  \begin{equation}
    \Vec{\mc{M}}(\F_{\rm C}) = \vv{OM}\pv\F_{\rm C}=\Vec{0}.
  \end{equation}
  \item \udl{TMC} : le moment cinétique de l'électron par rapport à l'origine $\rm O$ ou se situe le proton s'écrit pour une orbite circulaire.
  \begin{equation}
  \vv{L}_O=\vv{OM}\pv m_e\vv{v} = r \er \pv m_e v \et = m_e v r \ez.
  \end{equation}
  Le TMC indique alors que le moment cinétique est conservé,
  \eqbox[]{
  \vv{L}_O=\Vec{\cte}.
  }
\end{enumerate}


\parag{Moment magnétique orbital}
La trajectoire circulaire de l'électron, qui est une particule chargée, s'apparente à une boucle de courant d'intensité 
\begin{equation}
  I = \frac{-e}{T} = \frac{-e v}{2\pi r},
\end{equation}
où $T$ est la période du mouvement. Le moment magnétique de cette boucle de courant, appelé moment magnétique orbital, s'écrit
\begin{equation}
\M=IS\ez=\frac{-ev}{2\pi r}\pi r^2\ez=\frac{-evr}{2}\ez=\frac{-e}{2m_e}\vv{L}_O.
\end{equation}
On voit que le moment magnétique est proportionnel au moment cinétique orbital de l'électron.
Cette relation se généralise à toute particule de charge $q$.
% Cette relation se met sous la forme générale suivante, valable pour tout moment cinétique.
\begin{theorem}[Rapport gyromagnétique]\label{form:gyro}
  Une particule de masse $m$, de charge $q$, et de moment cinétique orbital $\Vec{L}$, a un moment magnétique
  $$
  \M=\gamma\vv{L},
  $$
  où $\gamma = \frac{q}{2m}$ est appelé \imp{rapport gyromagnétique}.
\end{theorem}
% Pour l'électron autour du noyau atomique, le rapport gyromagnétique vaut
% \begin{equation}
%   \gamma=\frac{-e}{2m}.
% \end{equation}

\begin{exemple}{Ordres de grandeur}
  On considère l'électron de l'atome d'hydrogène dans le modèle planétaire.
 Évaluer : le rapport gyromagnétique $\gamma$ de l'électron, la vitesse $v$ de l'électron, l'intensité du courant $I$ associé à la trajectoire de l'électron, et le moment magnétique orbital $\M$ de l'électron. 
 
 On donne $r = \SI{0.53}{\angstrom}$, $m_e = \SI{9.11e-31}{kg}$, et $e = \SI{1.60e-19}{C}$.

\end{exemple}
\examplesol{
  \begin{quest}
    \item	Le rapport gyromagnétique de l'électron vaut
    \eqbox{
    \gamma = \frac{-e}{2m_e} \simeq -\SI{8.8e10}{SI}.
    }
    \item La vitesse de l'électron s'obtient en appliquant la 2\ieme loi de Newton à l'électron, soit
    $$
    m_e \frac{v^2}{r} = \frac{e^2}{4\pi \epz r^2},
    $$
    d'où
    \eqbox{
    v = \sqrt{\frac{e^2}{4\pi \epz m_e r}} \simeq \SI{2.2e6}{m.s^{-1}}.
    }
    \item L'intensité du courant associé à la trajectoire de l'électron vaut
    \eqbox{
    I = \frac{-e}{T} = \frac{-e v}{2\pi r} \simeq \SI{0.40}{mA}.
    }
    \item Le moment magnétique orbital de l'électron vaut
    \eqbox{
    \M = I \pi r^2 = \frac{-evr}{2} \simeq \SI{2.0e-23}{A.m^2}.
    }
  \end{quest}
}

% \begin{remarque}
% (Hors programme). Le magnétisme dans la matière n'est en fait pas seulement dû au moment cinétique orbital. Les particules élémentaires se caractérisent également par une grandeur vectorielle appelée spin, qui n'a pas d'équivalent classique. Ce spin, équivalent à un moment cinétique intrinsèque $\vv{S}$, engendre également un moment magnétique de la forme $\M=\gamma\vv{S}$, mais la valeur du rapport gyromagnétique de spin diffère du rapport gyromagnétique orbital.
% \end{remarque}



\subsection{Magnétisme de spin}

\parag{L'expérience de Stern-Gerlach}

Cette expérience fut réalisé par O. \textsc{Stern} et W. \textsc{Gerlach} en 1922 dans le but de mesurer un moment magnétique atomique. Des atomes d'argent portés à l'état gazeux sont contenus dans une enceinte chauffée. Certains atomes sont éjectés de l'enceinte par une petite ouverture, formant ainsi un jet atomique dans la direction $(Ox)$. Ces atomes sont ensuite déviés par un champ magnétique transverse porté par $(Oz)$ produit par un électro-aimant. En sortie de la cavité, les atomes sont détectés sur un écran placé derrière la cavité. La forme des pièces polaires de l'aimant qui produit le champ magnétique est choisie de sorte à produire un gradient constant, de la forme $\B(M) = a z \ez$, avec $a>0$.

\begin{figure}
  \centering
  {\includegraphics[scale=1.1]{Stern-Gerlach.pdf}}
  \caption{Expérience de Stern-Gerlach. Un jet d'atomes d'argent est dévié par un champ magnétique non uniforme. Les atoms sont déviés vers le haut ou le bas suivant la composante $\mu_z$ du moment magnétique, puis projeté sur un écran. Le principe de l'expérience est décrit dans \href{https://fr.wikipedia.org/wiki/Fichier:Spin_quantique_et_exp\%C3\%A9rience_de_Stern_et_Gerlach.ogv}{cette vidéo}.}
  \label{fig:SG}
\end{figure}

Notons $\M$ le moment magnétique d'un atome d'argent et $\mu_z$ sa composante suivant l'axe du champ. L'atome étant de charge nulle, il ne subit aucune force de Lorentz. a force que subit un atome est celle exercée sur un dipôle magnétostatique, 
\begin{equation}
  \F = (\M \ps \grad) \B = \mu_z \dd{B}{z} \ez = \mu_z  a \ez.
\end{equation}
% \begin{remarque}
%   L'atom est neutre électriquement et donc il ne subit aucune force de Lorentz.
% \end{remarque}
La force est verticale et proportionnelle à $\mu_z$. On en déduit que la déviation de l'atome sur l'écran est également proportionnelle à $\mu_z$. Le comptage des points d'impact des atomes sur l'écran permet ainsi de mesurer la distribution de $\mu_z$ dans le jet d'atomes d'argent.

% \begin{exemple}{Calcul de la déviation dans l'expérience de Stern-Gerlach}
%   En appliquant le PFD, déterminer la hauteur d'impact $x$ des atomes sur l'écran en fonction.
%   %  de $a$, $D$, $v_0$, et $\M$.
%   \end{exemple}
%   \examplesol{
%     Le PFD appliqué à l'atome d'Ag de masse $m$ donne
%     \begin{equation}
%     \begin{cases}
%     m \dd{v_z}{t} & = 0 \\
%     m \dd{v_x}{t} & = m \mu_z a.
%     \end{cases}
%     \end{equation}
%     On note $v_z(\ell) = \cte = v_0$ la vitesse axiale de l'atome en entrée de la cavité. On a de plus $v_x(t) = \frac{\mu_z a}{m} t$. En supposant $v_x\ll v_z$, l'atome sort de la région soumis au champ magnétique à l'instant $t \simeq \ell/v_0$ et donc
%     \begin{equation}
%     v_x(\ell) = \frac{\mu_z a \ell}{m v_0}.
%     \end{equation}
%     On trouve un angle de déviation en sortie de cavité
%     \begin{equation}
%       \tan \alpha = \frac{v_x(\ell)}{v_z(\ell)} = \frac{ \mu_z a \ell}{m v_0^2}.
%     \end{equation}
%     L'atome se déplace ensuite en ligne droite et frappe l'écran à une hauteur
%     \eqbox{
%       z = D \tan \alpha  = \frac{ \mu_z a \ell}{m v_0^2}.
%     }  
%   }


\parag{Résultat et interprétation}
En physique classique, on s'attend à ce que la variable $\mu_z$ soit distribué \imp{uniformément} entre les valeurs extrémales $-\mu$ et $+\mu$, ce qui donnerait une distribution des atomes frappant l'écran plate quelque soit l'atome considéré, de la forme de la \Fig{SG}b). Or ce n'est pas ce qui est observé ! On observe au contraire deux taches symétriques correspondant bien séparées.

\begin{figure}[H]
  \centering
  \subfigure[]{\includegraphics[height=5cm]{SG_classique.pdf}}\hspace{0.5cm}
  \subfigure[]{\includegraphics[height=5cm]{SG_cadmium.pdf}}\hspace{0.5cm}
  \subfigure[]{\includegraphics[height=5cm]{SG_quantique.pdf}}\hspace{0.5cm}
  \subfigure[]{\includegraphics[height=5cm]{SG_postal.jpg}}
  \caption{ (a) Résultat attendu en physique classique. La distribution des points d'impact est plate pour tout atome. (b) Résultat obtenu pour des atomes de cadmium. La composante du moment magnétique est nul ! L'élargissement de la tache provient de la distribution des vitesses initiales des atomes. (c) Résultat obtenu pour des atomes d'argent. Les atomes se répartissent en deux taches symétriques correspondant : le moment cinétique est quantifié en deux valeurs opposées. (d) Résultat de l'expérience historique de Stern-Gerlach. Carte postale envoyée par S. \textsc{Gerlach} à N. \textsc{Bohr}.}
  \label{fig:SG}
\end{figure}

%  Les résultats de Stern et Gerlach conduisent à deux problèmes principaux.

\begin{enumerate}
    \item \udl{Règle de quantification}. 
    
    
    Imaginons que l'on réalise l'expérience de Stern-Gerlach avec des atomes de cadmium (ce n'est pas ce qu'ont réalisé expérimentalement Stern et Gerlach). On observerait alors que tous les atomes subissent une déviation nulle, soit $\mu_z=0$ pour le cadmium. D'après la \Form{gyro}, les électrons de valence de l'atome de cadmium ont donc un moment cinétique orbital nul, soit 
    \begin{equation}
      L_z = 0.
    \end{equation}
    On sait aujourd'hui que la configuration électronique de l'atome de cadmium est [Cd] = [Kr]4d$^{10}$5s$^2$. Les deux électrons de valence occupent une orbitale $s$, qui est donc de moment cinétique orbital nul. De façon générale, les orbitales $s$, $p$, $d$, ... sont associés à une nombre quantique orbital $l=0,1,2,...$, qui impose aux composantes du moment cinétique orbital suivant n'importe quel axe de prendre des valeurs entières de la constante de Planck réduite, $\hbar = \SI{1.05e-34}{J.s}$.

      \begin{prop}[Quantification du moment cinétique orbital]\label{prop:quant}
        Toute composante $L_z$  du moment cinétique \imp{orbital} est \imp{quantifiée} en multiple de $\hbar$, soit 
        \begin{equation*}
      L_z = m \hbar,
    \end{equation*}
    où $m=-\ell,-\ell+1,...,\ell$ pour une particule de nombre quantique orbital $\ell \in \N$.
  \end{prop}
  Pour des électrons de moment cinétique orbital $\ell$, on doit donc observer $2\ell+1$ de taches symétriques. 

  \begin{table}[h!]
    \centering
    \begin{tabular}{|c| c c c c c|}
      \hline
      Nombre quantique orbital $\ell$ & 0 & 1 & 2 & 3 & ... \\
      \hline
      Orbitale atomique & s & p & d & f & ... \\
      \hline
      Nombre quantique magnétique $m$ & $0$ & $-1,0,1$ & $-2,-1,0,1,2$ & $-3,-2,-1,0,1,2,3$ & ... \\
      \hline 
      Dégénérescence $2\ell+1$ & 1 & 3 & 5 & 7 & ...\\ 
      \hline
    \end{tabular}
    \caption{Correspondance entre le nombre quantique orbital $\ell$ et les orbitales atomiques.}
    \label{}
    
  \end{table}
  %   L'expérience montre la présence de deux tâches distinctes symétriques avec des atomes d'argent au lieu d'une seule tâche étendue. Cela montre que le moment magnétique $\M_z$, et donc le spin de l'électron $S_z$ ne peut pas prendre de valeurs continues mais uniquement deux valeurs discrètes opposées. Les mesures montres que ces valeurs de $S_z$ correspondent à $\pm\frac{\hbar}{2}$, où $\hbar$ est la constante de Planck réduite, avec $\hbar=\frac{h}{2\pi}$ et $h$ la constante de Planck qui vaut $h=\SI{6,62e-34}{J.s}$. Par analyse dimensionnelle, on remarque que $\SI{1}{J.s}=\SI{1}{kg.m^2.s^{-1}}$ qui est homogène ... à un moment cinétique! On dit que l'électron est une particule de spin $s=1/2$.
  %  Ce résultat se généralise à tout moment cinétique.

  \item \udl{Spin de l'électron}. 
  
  
  L'expérience historique montre que des atomes d'argent subissent une déviation répartie en deux taches symétriques. Or l'argent est de configuration électronique [Ag] = [Kr]4d$^{10}$5s$^1$. Le seul électron non apparié occupe la couche 5s, qui correspond à une orbitale de moment cinétique nul, soit $L_z = 0$, et donc aucune déviation ne devrait apparaître ! Qui plus est, un moment magnétique orbital ne peut pas conduire à une nombre \emph{paire} de taches.
  Ce résultat fut interprété par \textsc{Goudsmit} et \textsc{Uhlenbeck} en 1925, puis par W. \textsc{Pauli}, comme la preuve de l'existence d'un moment cinétique intrinsèque, de nature non orbitale, appelé \imp{spin}. 
  % On compris plus tard (W. Pauli) que les électrons (et toutes les particules élémentaires) possèdent un moment cinétique intrinsèque, présent même en l'absence de mouvement orbital. Ce moment cinétique est appelé \imp{spin} et est noté $\Vec{S}$. Il a même dimension que le moment cinétique orbital (en $\SI{}{kg.m^2.s^-1=J.s}$).
  
  \begin{prop}[Spin]\label{prop:quant}
    Toute particule possède un \imp{moment cinétique de spin} $\Vec{S}$ et un moment magnétique de spin
    \begin{equation*}
      \M_s = g\gamma\vv{S},
    \end{equation*}
    où $\gamma$ est le rapport gyromagnétique et $g$ est le \imp{facteur de Landé} ($g\simeq-2$ pour l'électron).
  \end{prop}
  Dans le cas d'un seul électron, on observe toujours deux taches symétriques, ce qui signifie que les composantes du spin sont quantifiées en demi-entiers de $\hbar$. On dit que l'électron est une particule de spin $s=1/2$. C'est aussi le cas de la plupart des particules élémentaires (fermions), comme les protons et les neutrons.
  \begin{prop}[Quantification du spin]\label{prop:}
    Toute composante du moment cinétique de \imp{spin} est \imp{quantifiée} en demi-entiers de $\hbar$, soit
    \begin{equation*}
  S_z = m_s \hbar,
\end{equation*}
où $m_s = -s, -s+1, ..., s$ pour une particule de spin $s$.
  \end{prop}
  Pour l'électron, les valeurs possibles du moment cinétique de spin  sont donc 
  \eqbox[]{
  S_z = \pm \frac{\hbar}{2}.
  }

  
\end{enumerate}

% \begin{figure}[H]
%   \centering
%   \includegraphics[width=0.5\textwidth]{SG_postal.jpg}
%   \caption{Résultat de l'expérience historique de Stern-Gerlach. Carte postale envoyée par Gerlach à Bohr.}
%   \label{fig:postal}
% \end{figure}


% De plus,  Mais alors pourquoi observe-t-on deux taches décentrées et pas une seule tache centrale ? On conpris donc plus tard avec W. Pauli que les électrons (et toutes les particules élémentaires constituant la matière) possèdent un moment cinétique intrinsèque, présent même en l'absence de mouvement orbital. Ce moment cinétique est appelé \imp{spin} et est noté $\Vec{S}$. L'électron est de spin $1/2$, ce qui signifie que les composantes du spin valent
% \begin{equation}
%   S_z =\pm \frac{\hbar}{2}.
% \end{equation}

% \begin{prop}[Spin]\label{prop:quant}
%   L'expérience de Stern-Gerlach montre que les électrons possèdent un moment cinétique intrinsèque, d'origine non orbitable, appelé \imp{spin}, noté $\Vec{S}$. Les composantes du spin sont \imp{quantifiées}, tout comme pour le moment cinétique orbital.
% \end{prop}



\parag{Magnéton de Bohr} L'unité de moment cinétique atomique naturelle est $\hbar$. Déterminons l'ordre de grandeur du moment magnétique d'un électron dans un atome. \begin{itemize}
  \item La projection du moment cinétique orbital $L_z$ d'un électron sur une couche $p$ vaut zéro, ou $L_z = \pm \hbar$. Le moment magnétique orbital de l'électron vaut alors $\mu_\ell = \gamma L_z \simeq \pm \frac{e\hbar}{2m_e}$.
  \item La projection du moment de spin $S_z$ d'un électron vaut $S_z = \pm \frac{\hbar}{2}$. Le moment magnétique de spin de l'électron vaut alors $\mu_s = -2\gamma S_z \simeq \pm \frac{e\hbar}{2m_e}$.
\end{itemize}
On en déduit que le moment magnétique atomique est de l'ordre de $\frac{e\hbar}{2m_e}$,  appelé magnéton de Bohr.
\begin{definition}[Magnéton de Bohr]\label{prop:}
  Le moment magnétique atomique est de l'ordre du \imp{magnéton de Bohr}
  \begin{equation*}
  \mu_{\rm{B}}=\frac{e\hbar}{2m_e} \simeq \SI{9.2e-24}{A.m^2},
 \end{equation*}
 avec $m_e$ la masse de l'électron et $e$ la charge élémentaire.
\end{definition}

\begin{exemple}{Ordre de grandeur de l'aimantation}
  Considérons un aimant permanent constitué d'atomes de fer en configuration $\alpha$, dont la configuration électronique est [Fe] = [Ar]3d$^6$4s$^2$. Les électrons appairés ne contribuent pas au moment cinétique orbital ou de spin. Seuls les deux électrons non appairés des orbitales $d$ contribuent au moment magnétique de spin du fer, de spin total $s=1$.
  Déterminer l'aimantation maximale de cet aimant.

  \textsc{Données} : masse molaire : $M({\rm Fe}) = \SI{56}{g.mol^{-1}}$ ;  masse volumique du fer $\alpha$ : $\rho({\rm Fe}) = \SI{7.9}{g.cm^{-3}}$.
  % Le moment magnétique de chaque atome de fer vaut alors $\M = -2\gamma S \hbar \simeq 4\mu_B$. Sachant que le fer possède environ $n=\SI{e29}{}$ atomes par mètre cube, on trouve une aimantation de l'ordre de $M = n \M \sim \SI{e5}{A.m^{-1}}$, ce qui correspond à l'aimantation typique d'un aimant permanent.
\end{exemple}
\examplesol{
  Le moment magnétique de spin ve chaque atome de fer vaut 
  $$
  \mu_s = -2\gamma S_z \simeq 2\mu_B \simeq \SI{1.8e-23}{A.m^2}.
  $$
  La densité volumique d'atomes de fer vaut
  $$
n = \frac{\rho({\rm Fe})}{M({\rm Fe})} \NA \simeq \SI{8.5e28}{m^{-3}}.
  $$
  L'aimantation maximale d'un aimant de fer s'obtient en supposant tous les moments magnétiques alignés, et vaut
  \eqbox{
  M = n \mu_s \simeq \SI{1.5e5}{A.m^{-1}}.
  }
}


\end{document}
